% BHU PYQ -- Manually transcribed & error-corrected
\documentclass[11pt,a4paper]{article}

\usepackage[a4paper,top=2.5cm,bottom=2.5cm,left=2.8cm,right=2.8cm,headheight=14pt]{geometry}
\usepackage{amsmath,amssymb}
\usepackage[T1]{fontenc}
\usepackage[utf8]{inputenc}
\usepackage{lmodern}
\usepackage{microtype}
\usepackage{fancyhdr}
\usepackage{enumitem}
\usepackage{hyperref}
\usepackage{lastpage}
\usepackage{parskip}

\hypersetup{colorlinks=true,urlcolor=black,linkcolor=black,
  pdftitle={MTB-507: Combinatorial Mathematics -- BHU PYQ},pdfauthor={Banaras Hindu University}}

\pagestyle{fancy}
\fancyhf{}
\renewcommand{\headrulewidth}{0.4pt}
\renewcommand{\footrulewidth}{0.4pt}
\fancyhead[L]{\small   \textit{Banaras Hindu University}}
\fancyhead[R]{\small   \textit{MTB-507: Combinatorial Mathematics}}
\fancyfoot[C]{\small Page  \thepage\ of \pageref{LastPage}}


\newlist{parts}{enumerate}{2}
\setlist[parts,1]{label=(\alph*),leftmargin=*,itemsep=5pt,topsep=3pt}
\setlist[parts,2]{label=(\roman*),leftmargin=*,itemsep=3pt,topsep=2pt}

\newcommand{\pts}[1]{\hfill   \textit{[#1]}}


\begin{document}


\begin{center}
  {\large\bfseries BANARAS HINDU UNIVERSITY}\\[2pt]
  {\normalsize B.A./B.Sc. (Hons.) Semester V Examination, 2022-23}\\[6pt]
  \rule{0.8\linewidth}{0.5pt}\\[6pt]
  {\large\bfseries Paper No. MTB-507: Combinatorial Mathematics}\\[4pt]
  {\normalsize Time: 3 Hours\hfill Maximum Marks: 70}
\end{center}

\rule{\linewidth}{0.4pt}

\smallskip



\noindent   \textit{Instructions: Answer   \textbf{five} questions including Question~1 which is compulsory. Symbols have their usual meanings.}

\medskip


\noindent   \textbf{Question 1.}

\begin{parts}
  \item In how many ways can four tyres be put on a car? \pts{1}
  \item A king is placed on the bottom left-hand square of an $8  \times 8$ chessboard and is to move to the top right-hand corner square. If it can move only up or to the right, how many possible paths does it have to choose from? \pts{2}
  \item It is required to seat $n$ people round a table. Show that this can be done in $(n-1)!$ ways. \pts{3}
  \item A pack of 52 cards is divided among 4 people so that each gets 13 cards. How many such deals are possible? \pts{2}
  \item If
  \[ B = 
\begin{pmatrix} 1 & 1 & 1 & 1 \\ 1 & -1 & 1 & -1 \\ 1 & 1 & -1 & -1 \\ 1 & -1 & -1 & 1 \end{pmatrix} \]
  derive from $B$ an $8  \times 8$ Hadamard matrix and hence write down the incidence matrix of a $(7,3,1)$-configuration. \pts{3}
  \item Find the Rook's polynomial of the following board:
  \[ 
\begin{array}{ccc}
  \square & \square & lacksquare \\
  \square & lacksquare & \square \\
  lacksquare & \square & \square
  \end{array} \]
  where $lacksquare$ represents a shaded/unavailable square. \pts{3}
\end{parts}

\medskip


\noindent   \textbf{Question 2.}

\begin{parts}
  \item Suppose that each of $k$ indistinguishable golf balls has to be coloured with any one of $n$ given colors and we have to find the different colourings. Show that the number of different colourings is same as the number of non-negative integral solutions of the equation
  \[ x_1 + x_2 + \cdots + x_n = k \]
  Let $f(n,k)$ be the number of solutions of this equation. Show that:
  
\begin{parts}
    \item $f(n,k) = f(n-1,k) + f(n,k-1)$ for all $n \ge 2, k \ge 2$.
    \item $f(n,k)$ is the coefficient of $x^k$ in $(1-x)^{-n}$.
    \item $f(n,k)$ is the number of ways of choosing $(n-1)$ places out of $(n+k-1)$ places.
  \end{parts} \pts{10}
  \item In a cricket league of $n$ teams, each team plays each other twice. The number of games played is therefore $2C$, where $C$ is the number of ways of choosing 2 objects from $n$ given objects. Prove that
  \[ 2C = 2\left[ (n-1)+(n-2)+\cdots+2+1 \right] \]
  and deduce the number of games played in a league of 22 teams. \pts{4}
\end{parts}

\medskip


\noindent   \textbf{Question 3.}

\begin{parts}
  \item By using the identity
  \[ (1+x)^{2n} = (1+x)^n(1+x)^n \]
  and considering the coefficient of $x^n$ on both sides, prove that
  \[ inom{2n}{n} = inom{n}{0}^2 + inom{n}{1}^2 + inom{n}{2}^2 + \cdots + inom{n}{n}^2 \]
  Verify this in the case $n=5$. \pts{7}
  \item Use the binomial theorem to evaluate:
  
\begin{parts}
    \item $\sum_{r=0}^{n} (-1)^r inom{n}{r}$
    \item $\sum_{r=0}^{n} r inom{n}{r}$
  \end{parts} \pts{7}
\end{parts}

\medskip


\noindent   \textbf{Question 4.}

\begin{parts}
  \item Show that, if a graph has $2n$ vertices, each of degree $\ge n$, then the graph has a perfect matching. Show by an example that the converse of this statement may not be true. \pts{6}
  \item State and prove Hall's marriage theorem. \pts{8}
\end{parts}

\medskip


\noindent   \textbf{Question 5.}

\begin{parts}
  \item Let $\{a_n\}$ be a sequence satisfying $a_n = A a_{n-1} + B a_{n-2}$, $n \ge 2$, where $a_0$ and $a_1$ are given. $A$ and $B$ are constants. If $r_1$ and $r_2$ are the roots of $r^2 - Ar - B = 0$, then show that:
  
\begin{parts}
    \item $a_n = k_1 r_1^n + k_2 r_2^n$, if $r_1$ and $r_2$ are real and distinct;
    \item $a_n = (k_1 + n k_2)r^n$, if $r_1 = r_2 = r$;
  \end{parts}
  where $k_1, k_2$ are constants uniquely determined by the initial conditions. \pts{7}
  \item Solve the Fibonacci recurrence relation: $F_n = F_{n-1} + F_{n-2}$ for $n \ge 2$, with $F_0 = 0$, $F_1 = 1$. \pts{7}
\end{parts}

\medskip


\noindent   \textbf{Question 6.}

\begin{parts}
  \item Define derangements. If $D_n$ denotes the number of possible derangements of $n$ symbols, then find $D_n$. \pts{7}
  \item Given a chessboard $C$, choose any square of $C$ and let $D$ denote the board obtained by deleting from $C$ every square in the same row or column as the chosen square (including the chosen square itself). Let $E$ denote the board obtained from $C$ by deleting only the chosen square. Then show that:
  \[ R(x, C) = x R(x, D) + R(x, E) \] \pts{7}
\end{parts}

\medskip


\noindent   \textbf{Question 7.}

\begin{parts}
  \item For a $(v,b,r,k,\lambda)$-configuration, show that $vr = bk$, $r(k-1) = \lambda(v-1)$, and $b \ge v$. \pts{10}
  \item Show that:
  
\begin{parts}
    \item A code will detect all sets of $h$ or fewer errors if any two words differ in at least $(h+1)$ places. \pts{4}
    \item A code will correct all sets of $h$ or fewer errors if any two words differ in at least $(2h+1)$ places. \pts{4}
  \end{parts}
\end{parts}



\vfill
\rule{\linewidth}{0.4pt}

\begin{center}\small
\href{https://bhu-exam.vercel.app/}{Click here to visit our website}\\[4pt]\href{https://me-aryan.vercel.app/}{Click here to visit our contributor}\\[4pt]\href{https://quashan-me.vercel.app/}{Click here to visit our contributor}
\end{center}

\end{document}
